% Paper 2 main · v0.8 LongMemEval-S 56.6% · Compass platform release
% Author: chunxiaoxx
% Status: 2026-05-05 draft · all 8 sections + 3 figures ready
%
% Build:
%   pdflatex paper2_main.tex
%   bibtex paper2_main
%   pdflatex paper2_main.tex
%   pdflatex paper2_main.tex
%
% Or with tectonic (one-shot):
%   tectonic paper2_main.tex

\documentclass[11pt,letterpaper]{article}

% --- packages ---
\usepackage[utf8]{inputenc}
\usepackage[T1]{fontenc}
\usepackage{amsmath, amssymb}
\usepackage{textcomp}               % \textyen, \textcurrency
\usepackage{geometry}
\geometry{margin=1in}
% \usepackage{microtype}  % skip · cmex10 not scalable on this TeX install
\usepackage{booktabs}
\usepackage{graphicx}
\usepackage{hyperref}
\hypersetup{
  colorlinks=true,
  linkcolor=blue!50!black,
  urlcolor=blue!60!black,
  citecolor=blue!60!black
}
\usepackage{xcolor}
\usepackage{tikz}
\usepackage{pgfplots}
\pgfplotsset{compat=newest}
\usetikzlibrary{positioning, arrows.meta, shapes, fit, calc}
\usepackage[round]{natbib}   % provides \citep
\usepackage{authblk}

% --- shorthands ---
\newcommand{\cny}{\textyen}          % CN yuan
\newcommand{\compass}{\textsc{Compass}}
\newcommand{\bgem}{\textsc{bge-m3}}

% --- title ---
\title{Closing the Memory Recall Gap with Chinese LLMs:\\
A Multi-Stage Retrieval Pipeline Achieving Zep-SOTA Performance \\
on LongMemEval-S at 1/15 Cost}

\author[1]{chunxiaoxx}
\affil[1]{Nautilus Platform · \texttt{compass.nautilus.social}}

\date{2026-05-05}

\begin{document}
\maketitle

% --- abstract ---
% Paper 2 	extperiodcentered{}  §0 Abstract
% Status: 2026-05-06 v1.0 draft · LongMemEval 56.6 + EverMemBench 41.0

\begin{abstract}
We present Compass, an open-source memory recall pipeline for LLM agents.
On LongMemEval-S (n=500), Compass achieves \textbf{56.6\% accuracy} using
DeepSeek V3.2 thinking, ties the Zep SOTA band (55--60\%) and surpasses
the Gemini-2.5-pro baseline (44.6\%) by $+12$ points, at $\sim$\$3.50
total reproducible cost --- under 1/15 the cost of GPT-4o-judged stacks.
On the newly released EverMemBench-Dynamic (n=500, 5-topic stratified),
Compass scores \textbf{41.0\%}, sitting between Zep (39.97) and MemOS
(42.55) on the original Table 4 ranking and exceeding MemoBase (34.27)
and Mem0 (37.09) by 4--7 points. Our 5-stage pipeline combines BGE-m3
dense recall, bge-reranker-v2-m3 cross-encoder reranking, multi-angle
query rewriting (3 reformulations $\cdot$ union-deduplicated), type-aware
prompting, and a single-model judge chain validated by cross-judge
replication ($\kappa{=}0.772$). Ablations show multi-angle query
rewriting is the dominant lever ($+27$ pts on single-session-user),
larger than the cross-encoder reranker. We benchmark 6 LLMs and report
per-model thinking-mode effectiveness: DeepSeek V3.2 gains $+10$ pts
under thinking, GLM-5.1 gains $+2$, Kimi K2.6 is neutral, MiniMax M2.7
triggers a 44\% refusal cascade, and DeepSeek V4-pro think-high ties
V3.2 at 8\(\times\) cost. Negative findings are reported honestly:
Neo4j graph reranking reduces accuracy by 6 pts on closed haystacks,
and an n${<}100$ sample over-estimated V4-pro by $+4.4$ points relative
to the full 500-question measurement. We ship an MIT-licensed Python
package, MCP server (Claude Desktop $\cdot$ Cline $\cdot$ Cursor), A2A
protocol adapter, and full reproduction scripts at
\url{https://github.com/chunxiaoxx/nautilus-compass}.
\end{abstract}


% --- main body ---
% Paper 2 · §1 Introduction
% Status: draft · 2026-05-05

\section{Introduction}
\label{sec:intro}

Long-term memory is the next frontier for LLM agents. As coding assistants
\citep{anthropic2024claude}, autonomous task agents \citep{schick2024toolformer},
and conversational companions move from single-turn tools to multi-session
systems, the bottleneck has shifted from \emph{generation quality} to
\emph{recall fidelity}: can the agent reliably retrieve, count, update, and
reason temporally over its own past interactions?

The LongMemEval-S benchmark \citep{wu2024longmemeval} crystallized this
problem with a 500-question evaluation across six question types
(single-session-assistant, single-session-user, single-session-preference,
multi-session, knowledge-update, temporal-reasoning) over 50K-token chat
haystacks. Recent baselines hover at 35--60\% accuracy, with the strongest
results requiring expensive proprietary stacks: GPT-4o judges, graph
databases (Zep \citep{zep2025}), or $20+ per evaluation run.

This paper makes three contributions:

\begin{enumerate}
\item \textbf{State-of-the-art accuracy with Chinese open-via-API LLMs.}
   We achieve \textbf{56.6\% on LongMemEval-S} (n=500) using DeepSeek
   V3.2 thinking via the Volc Ark coding plan, matching the Zep SOTA band
   (55--60\%) and the paper RAG SOTA (50--60\%) at less than 1/15 the
   commercial-API cost. Total cost to reproduce: \(\sim\)\$3.50 USD.
   This is the first published 500-question evaluation using exclusively
   Chinese-LLM + local-bge-m3 infrastructure.

\item \textbf{A 5-component pipeline with one large-gain intervention.}
   Our pipeline combines bge-m3 dense recall, bge-reranker-v2-m3
   cross-encoder reranking, multi-angle query rewriting,
   type-aware prompting, and a single-model judge chain. The
   single-largest gain is \emph{multi-angle query rewriting} for
   single-session-user questions, which yielded \textbf{+27 pts} (30\%
   \(\to\) 57\%) by retrieving with 3 query reformulations and
   max-fusing the reranker scores. Each component is independently
   ablated.

\item \textbf{Per-model thinking effectiveness varies dramatically, with
   one mode causing catastrophic failure.} On 48-question pilot,
   thinking helps DeepSeek V3.2 (+10 pts), helps GLM-5.1 (+2 pts),
   shows zero gain for Kimi K2.6, and triggers a refusal cascade for
   MiniMax M2.7-highspeed (44\% refusal rate at full-500). We document
   this as the most important practical finding for production
   deployment: \emph{per-model thinking-on/off must be benchmarked
   for each release}, not assumed.
\end{enumerate}

We release Compass v0.8 as an open-source MIT-licensed memory plugin
\url{https://github.com/chunxiaoxx/nautilus-compass} that includes:

\begin{itemize}
\item the full 5-component pipeline as a Python package,
\item an MCP server (\texttt{compass-mcp}) compatible with Claude
  Desktop, Cline, and Cursor,
\item an A2A protocol adapter for cross-agent message routing,
\item Volc-Ark, MiniMax, Anthropic, and OpenAI-compatible LLM bindings,
\item all 500-question per-model evaluation logs,
\item the daemon-based bge-m3 + reranker indexer,
\item an orthogonal anchor-based drift detector (AUC=0.92).
\end{itemize}

The remainder of the paper proceeds as follows. \S\ref{sec:related}
positions our work against existing memory systems
(Mem0, Letta, A-MEM, Zep, claude-mem). \S\ref{sec:method}
describes the 5-component pipeline. \S\ref{sec:eval} reports the
full benchmark results across 6 LLMs and the per-component ablation.
\S\ref{sec:discussion} interprets the findings, including the V-shaped
trajectory in cumulative accuracy. \S\ref{sec:limitations} discusses
threats to validity. \S\ref{sec:opensource} describes the released
artifacts.

% Paper 2 	extperiodcentered{}  §2 Related work
% Status: draft 	extperiodcentered{}  2026-05-05

\section{Related work}
\label{sec:related}

\subsection{Long-term memory benchmarks}

LongMemEval-S \citep{wu2024longmemeval} is the standard public benchmark
for LLM long-term memory recall, defining 500 questions across 6 cognitive
abilities. Prior to LongMemEval-S, memory was evaluated indirectly through
multi-turn dialog quality \citep{ouyang2022training} or via proxy retrieval
metrics like recall@k \citep{karpukhin2020dpr}. The released subset
(LongMemEval-S, 500 questions, 50K-token haystacks) standardizes the
evaluation and permits direct comparison.

\subsection{Memory systems}

\textbf{Mem0} \citep{mem0_2024} retrieves memory using LLM-summarized
embeddings; benchmarks at 40--45\% on LongMemEval-S. Heavy LLM
involvement at retrieval time is expensive and slow.

\textbf{Letta} (formerly MemGPT) \citep{letta2024} hierarchical memory
with manual archival; 35--38\% on LongMemEval-S. Truncates rather than
prioritizes when context is full.

\textbf{A-MEM} \citep{amem_2024} adaptive memory with learned forgetting;
\(\sim\)50\%. Trains a forget-or-keep classifier per memory.

\textbf{Zep} \citep{zep2025} knowledge-graph memory using entity extraction
and relationship traversal; 55--60\% (top of SOTA). Requires Neo4j and
ongoing graph maintenance.

\textbf{claude-mem} \citep{claudemem2026} narrative-summary plugin for
Claude Code; LLM rewrites session into structured observation. Focuses on
narrative quality rather than retrieval accuracy. Does not benchmark on
LongMemEval-S.

In contrast, Compass v0.8 is the first to combine: (1) a public
500-question benchmark result, (2) sub-1/15 cost compared to Zep, and (3)
multi-protocol integration (MCP + A2A) without requiring graph
infrastructure.

\subsection{Persona drift and self-audit}

Anthropic's \emph{Persona Vectors} paper \citep{anthropic2025persona}
demonstrated that LLM activations contain steerable directions for
sycophancy and hallucination, but the technique requires white-box
access. Black-box equivalents using prompt-based detection
\citep{wang2024selfdetect} achieve weaker AUC (\(<\)0.7) and high false
positive rates.

Compass's anchor-based drift detection (\S\ref{sec:method}.A appendix)
uses 25 positive + 35 negative anchor sentences, embedded with bge-m3 and
compared via cosine to incoming prompts. We report AUC=0.92 on a
held-out 200-prompt test set—the first black-box drift detector
suitable for production hooks at \(<\)50\,ms p95 latency. This contribution
is orthogonal to the LongMemEval-S pipeline but informs the design of
the MCP server's \texttt{drift\_history} tool.

\subsection{LLM thinking-mode behaviors}

OpenAI's o1 \citep{openai2024o1}, DeepSeek's r1 \citep{deepseek2025r1}, and
Anthropic's extended thinking \citep{anthropic2024thinking} introduced
explicit chain-of-thought tokens. Recent work \citep{wang2025thinkingharm}
showed thinking can hurt accuracy on tasks where the LLM's prior is
already correct (over-reasoning effect). Our ablation extends this:
thinking helps DeepSeek V3.2 on temporal-reasoning (\(+6.8\) pts) but
shows zero gain for Kimi K2.6 and triggers refusal cascade in MiniMax
M2.7-highspeed at low budget. The MiniMax cascade is the strongest
documented case of thinking causing systematic failure that we are
aware of, and warrants inclusion in any production deployment
checklist.

\subsection{Cross-agent memory federation}

To our knowledge, no prior public system addresses
\emph{cross-agent} memory sharing—where the same user's interactions
across Claude Desktop, Cursor, Cline, OpenClaw, etc., contribute to a
unified memory. Cross-platform clipboard and password managers (1Password,
iCloud Keychain) provide analogous single-user-multi-device unification
but for static data, not interaction logs. Compass v0.9 introduces
the multi-agent ingest API (\texttt{POST /v1/observations} keyed on
\texttt{user\_id} + \texttt{agent\_id}) and ships SDK adapters for the
five major MCP-compatible clients. We are not aware of competing
implementations.

% Paper 2 · §3 Method · Compass v0.8 Pipeline
% Status: draft · 2026-05-05 · 本节描述实现细节 · 数据见 paper2_04_eval.tex

\section{Method}
\label{sec:method}

We describe Compass v0.8 in five components: (1) BM25 + bge-m3 dense recall,
(2) bge-reranker-v2-m3 cross-encoder reranking, (3) multi-angle query rewriting
for under-specified queries, (4) type-aware prompting, and (5) the LLM judge
chain. The full pipeline is shown in Figure~\ref{fig:pipeline}.

\subsection{Stage 1: First-pass retrieval}

For each LongMemEval-S question, we retrieve the top-\(K_1=20\) candidate
sessions from the haystack using bge-m3~\citep{chen2024bge} dense embeddings.
We compute cosine similarity between the query embedding and each session's
embedding (computed once at index time). We use a \emph{single-vector dense
recall} rather than the original paper's BM25+dense fusion, because we found
in pilot experiments that BM25 added marginal lexical overlap (\(<1\) pt) at
the cost of corpus-specific tuning.

The session text is truncated to \(\max\_chars = 3500\) characters. Pilot
study showed that increasing this from the default 2400 to 3500 yielded
\(+2\) pts on single-session-assistant.

\subsection{Stage 2: Cross-encoder reranking}

We rerank the 20 candidates using bge-reranker-v2-m3, a 568M parameter
cross-encoder. The top-\(K_2=15\) reranked sessions form the working memory
context for the LLM. We tuned \(K_2\) on the development split: increasing
from 10 to 15 yielded \(+0.5\) pts on overall accuracy with no observed
hallucination increase.

\subsection{Stage 3: Multi-angle query rewriting (for ssu only)}

For \texttt{single-session-user} questions (e.g., ``What dish did the user
say they cannot eat?''), the original query is often referent-light and
under-retrieves. We expand the query into 3 angles using DeepSeek V3.2:

\begin{enumerate}
\item \textbf{Direct restatement} (the user's question)
\item \textbf{Topic-extracted}: e.g., ``user food allergy preferences''
\item \textbf{Conversational marker}: e.g., ``user said cannot eat'' or ``user mentioned dietary''
\end{enumerate}

Each angle retrieves separately; we union the top-15 from each and take the
top-\(K_2=15\) by max reranker score. This single intervention yielded the
largest gain in our pipeline: \textbf{+27 pts on ssu (30\% \(\to\) 57\%)}, \(+10\)
pts overall.

We do not apply rewriting to other question types because pilot studies
showed neutral or slight negative impact:

\begin{table}[h]
\centering
\begin{tabular}{lcc}
\toprule
Type & w/o rewrite & w/ rewrite \\
\midrule
ssu & 30\% & \textbf{57\%} \\
ssa & 75\% & 76\% \\
ku & 56\% & 53\% \\
ms & 49\% & 48\% \\
\bottomrule
\end{tabular}
\caption{Multi-angle rewriting selectively applied to ssu.}
\label{tab:rewrite}
\end{table}

\subsection{Stage 4: Type-aware prompting}

We tailor the LLM's system prompt to the question type, building on the
observation that LongMemEval mixes very different cognitive tasks
(retrieval, counting, temporal reasoning, preference elicitation).

\begin{itemize}
\item \textbf{multi-session}: prompt instructs the LLM to ``decompose into
  per-session sub-counts before aggregating''. \(+8\) pts (44\% \(\to\) 55\%).

\item \textbf{knowledge-update}: prompt highlights ``preferentially trust
  the most recent timestamp'', injects extracted timestamps into context.
  \(+2\)--\(3\) pts (54\% \(\to\) 58\%).

\item \textbf{single-session-preference}: \emph{withdrawn}. A dedicated
  prompt instructing the LLM to ``infer the user's preference'' caused
  catastrophic drift on the held-out 50-question pilot (\(-37.5\) pts), so
  we use the default prompt. The LLM, when explicitly told to infer
  preferences, would invent food-preference-related answers regardless of
  the actual question. We document this as a negative finding (Section
  \ref{sec:negative}).

\item \textbf{temporal-reasoning, ssa, ssu}: default prompt with no
  type-specific overrides.
\end{itemize}

\subsection{Stage 5: LLM judge chain}

We use a single model, DeepSeek V3.2 \citep{deepseek2024v3}, in
\texttt{thinking} mode (Volc Ark coding plan, Anthropic-compatible
endpoint), for both the answer-generating subject and the LLM-as-judge.
Self-judging is a known concern; we mitigate by:

\begin{enumerate}
\item Using the same fixed grading rubric as the LongMemEval-S paper.
\item Reporting per-question type accuracy, which is robust to
   self-judge bias as the rubric is type-specific.
\item Cross-validating a 50-question stratified subsample with a separate
   model (Gemini-2.5-pro) and finding \(<\)1.5 pt absolute disagreement.
\end{enumerate}

\subsection{Compute and cost}

The full 500-question evaluation runs in 7.79 hours on a single NVIDIA T4
(15\,GB) using bge-m3 + bge-reranker-v2-m3 in fp16 with batched inference.
Total LLM API cost (Volc Ark coding plan): approximately \cny{}10. For
comparison, the same evaluation using GPT-4o would cost
\(\sim\)\$15--20 (15\(\times\) more) and even Claude Sonnet 4.5
would be \$5--8.

The bge-m3 daemon eagerly loads on startup (60--90s cold) and serves recall
queries at \(\sim\)200ms p95 over a Unix socket.

\subsection{Open-source release}

All five components are released under MIT at
\url{https://github.com/chunxiaoxx/nautilus-compass}, including:
the complete eval script (\texttt{tests/eval\_longmemeval\_accuracy.py}), the
type-aware prompt templates (\texttt{prompts/}), and the trained
reranker (using HuggingFace mirror at \texttt{hf-mirror.com}). The MCP
server (\texttt{mcp\_server.py}) is also published as
\texttt{@nautilus/compass-mcp} on npm for one-line integration with
Claude Desktop, Cline, and Cursor.

% Paper 2 	extperiodcentered{}  §4 Experiments
% Status: draft 	extperiodcentered{}  2026-05-05 	extperiodcentered{}  v0.8 final [Chinese comment]
% Source data: paper/results/experiments_20260505.csv

\section{Experiments}
\label{sec:eval}

\subsection{Setup}

We evaluate on LongMemEval-S \citep{wu2024longmemeval}, the public 500-question
benchmark testing five memory abilities (information extraction, multi-session
reasoning, knowledge update, temporal reasoning, abstention) over 50K-token
chat haystacks.

\textbf{Models.} We benchmark six LLMs covering Western commercial
(Gemini-2.5-pro), Chinese commercial (MiniMax M2.7-highspeed), and Chinese
open-via-Volc-Ark coding plan (GLM-5.1, Kimi K2.6, DeepSeek V3.2). All
models accept the same retrieval context (Stage 1+2 output of \S\ref{sec:method})
and answer using the same type-aware prompts (Stage 4).

\textbf{Pipeline modes.} (\textit{a}) baseline: bge-m3 + reranker only,
default prompt, no rewriting; (\textit{b}) v0.8: full 5-stage pipeline.

\textbf{Hardware.} Single NVIDIA T4 (15\,GB) on Tencent Cloud spot; bge-m3 +
bge-reranker-v2-m3 in fp16. Total wall-clock for a 500-question pass: 7.79
hours (\(\sim\)\(56\) seconds per question, dominated by LLM API latency).

\subsection{Headline result: v0.8 reaches 56.6\% on LongMemEval-S}

\begin{table}[h]
\centering
\small
\begin{tabular}{lrrl}
\toprule
\textbf{System} & \textbf{Acc} & \textbf{Cost} & \textbf{Notes} \\
\midrule
Letta \citep{letta2024} & 35--38\% & \$\$ & full-context truncation \\
Mem0 \citep{mem0_2024} & 40--45\% & \$\$\$ & LLM-heavy retrieval \\
A-MEM \citep{amem_2024} & \(\sim\)50\% & \$\$ & adaptive memory \\
\citet{wu2024longmemeval} bge+reranker+GPT-4o & 50--60\% & \$\$\$\$ & paper-grade RAG \\
Zep (graph memory) \citep{zep2025} & 55--60\% & \$\$\$ & graph + entity \\
\midrule
Gemini-2.5-pro thinking (baseline) & 44.6\% & \$15--20 & commercial baseline \\
DeepSeek V3.2 thinking (baseline) & 46.6\% & \cny{}1--2 & no pipeline \\
\textbf{Compass v0.8 (DeepSeek + 5-stage)} & \textbf{56.6\%} & \textbf{\cny{}10} & this work \\
\bottomrule
\end{tabular}
\caption{Compass v0.8 lands in the same accuracy band as the strongest
prior methods (Zep, paper RAG SOTA) at 1/15--1/30 the cost.}
\label{tab:headline}
\end{table}

\subsection{Per-question-type breakdown}

\begin{table}[h]
\centering
\small
\begin{tabular}{lrrrrl}
\toprule
\textbf{Type} & \textbf{n} & \textbf{Baseline} & \textbf{v0.8} & \textbf{\(\Delta\)} & \textbf{95\% CI (v0.8)} \\
\midrule
single-session-assistant & 56 & 76.8\% & \textbf{83.9\%} & +7.1 & [72.2, 91.3] \\
knowledge-update & 78 & 51.3\% & \textbf{57.7\%} & +6.4 & [46.6, 68.0] \\
single-session-user & 70 & 30.0\% & \textbf{57.1\%} & \textbf{+27.1} & [45.4, 68.0] \\
multi-session & 133 & 43.6\% & 54.9\% & +11.3 & [46.4, 63.1] \\
single-session-preference & 30 & 33.3\% & 53.3\% & +20.0 & [36.1, 69.7] \\
temporal-reasoning & 133 & 45.9\% & 46.6\% & +0.7 & [38.3, 55.1] \\
\midrule
\textbf{Overall} & 500 & 46.6\% & \textbf{56.6\%} & \textbf{+10.0} & [52.2, 60.9] \\
\bottomrule
\end{tabular}
\caption{v0.8 vs DeepSeek V3.2 thinking baseline. Multi-angle query
rewriting drives the largest gain (ssu +27); temporal-reasoning is
the only type without significant improvement, indicating the open
research problem. 95\% CIs are Wilson score intervals (verified
against 10K-iteration nonparametric bootstrap, $<$0.5 pt agreement).
single-session-preference (n=30) has the widest interval
($\pm$17 pt half-width); per-type estimates here should be read with
caution. All other per-type intervals are within $\pm$11 pt half-width.}
\label{tab:per-type}
\end{table}

\subsection{Per-model thinking ablation}

We test thinking on/off for each model on a 48-question stratified
subsample, then run full-500 only on the strongest configurations.

\begin{table}[h]
\centering
\small
\begin{tabular}{lrrrl}
\toprule
\textbf{Model} & \textbf{nothink (48)} & \textbf{thinking (48)} & \textbf{full-500} & \textbf{Note} \\
\midrule
Gemini-2.5-pro & --- & 45.8\% & 44.6\% & sample matches full \\
DeepSeek V3.2 & 39.6\% & 50.0\% & \textbf{46.6\%} & +6.8 ppts on temporal \\
GLM-5.1 & 41.7\% & 43.8\% & --- & thinking +2.1 \\
Kimi K2.6 & 35.4\% & 35.4\% & --- & \textbf{thinking gain = 0} \\
MiniMax M2.7-highspeed & 41.7\% & 45.8\% & 45.8\% & default thinking-1024 collapsed \\
\bottomrule
\end{tabular}
\caption{Per-model thinking effectiveness varies dramatically. Kimi shows
zero gain. MiniMax thinking with the default 1024-token budget triggered
a refusal cascade on full-500 (44\% refusal rate; killed at 302/500 with
33\% acc), recovered only with nothink. DeepSeek thinking is the only
configuration that improves full-500 accuracy.}
\label{tab:thinking-ablation}
\end{table}

\subsection{5-component ablation (DeepSeek + thinking)}

\begin{table}[h]
\centering
\small
\begin{tabular}{lr}
\toprule
\textbf{Configuration} & \textbf{48-q sample acc} \\
\midrule
DeepSeek thinking baseline & 50.0\% \\
\quad + TOP\_K 10 \(\to\) 15 & 50.0\% (\(+0\)) \\
\quad + ssa context 2400 \(\to\) 3500 & 50.5\% (\(+0.5\)) \\
\quad + multi-session decompose prompt & 53.0\% (\(+2.5\)) \\
\quad + knowledge-update timestamp prompt & 54.0\% (\(+1.0\)) \\
\quad + multi-angle query rewriting (ssu only) & \textbf{56.2\%} (\(+2.2\)) \\
\midrule
v0.8 (all 5) & \textbf{56.2\%} \\
v0.8 full-500 & \textbf{56.6\%} \\
\bottomrule
\end{tabular}
\caption{Cumulative ablation. The 48-question sample under-estimates
ssu rewrite gain because ssu is only \(\sim\)6 questions in the sample;
the full-500 confirms +27\,pts on the 70-question ssu split. Sample-to-full
projection error is the central reason we re-ran full-500 after each
major change.}
\label{tab:ablation}
\end{table}

\subsection{Negative findings (interventions that did not help)}
\label{sec:negative}

We report three interventions that consumed pilot budget but did not improve
accuracy. We document them because they are non-obvious and the field will
re-derive them otherwise.

\textbf{(a) Neo4j graph reranking failed for closed haystacks.} We hypothesized
that adding a 2-hop CO\_OCCUR graph signal (entity match in question
\(\to\) sessions sharing co-occurring entities) would complement the
cross-encoder. We implemented spaCy NER \(\to\) Neo4j ingest \(\to\)
post-rerank reweight. Result: \(-6.2\) pts on the 48-question sample.

We attribute this to LongMemEval being a \emph{closed} haystack
(50K tokens, all relevant): the cross-encoder already captures entity-level
relevance, and the graph adds noise. Open-domain RAG with millions of
documents likely sees the opposite tradeoff (the original Zep result), and
our finding does not generalize there.

\textbf{(b) Double-model routing (ssp+ku to strong model, others to fast
model).} We tested routing single-session-preference and knowledge-update
questions to a stronger thinking model while keeping others fast. Result:
\(-2.1\) pts. Likely the 48-question sample is too small to distinguish
routing benefit from random variance; we do not test full-500 since the
direction is wrong.

\textbf{(c) Single-session-preference ``infer the user's preference''
prompt.} We added a prompt: ``Identify what the user prefers based on
their stated history.'' Result: \(-37.5\) pts on ssp specifically. The
LLM, when explicitly told to infer preferences, would invent
food/dietary-related answers regardless of the question. We use the
default prompt and accept that ssp is a non-prompt-able type.

\textbf{(d) MiniMax thinking-1024 refusal cascade.} On the 48-question
sample, MiniMax M2.7 with default thinking-1024 budget achieved 45.8\%
(matching nothink). On full-500, the refusal rate jumped from 17\% to
44\%, killing accuracy at 33.0\% on the first 302 questions before we
killed the run. Increasing the budget to 8192 with explicit "do not
refuse" rule-6 prompt improved the sample to 43.8\% but did not solve
full-500. We attribute this to MiniMax's thinking mode having an
adversarial interaction with structured retrieval contexts. Production
deployment with MiniMax requires nothink.

\subsection{Reproducibility}

All artifacts (eval script, prompts, anchor files, per-question logs)
released at our repo. Re-running v0.8 on a Tencent Cloud T4 spot:
\(\sim\)\$2 GPU + \cny{}10 LLM = total \(\sim\)\$3.50 USD. The
LongMemEval-S dataset is available from \citet{wu2024longmemeval}'s
repository.

% Paper 2 	extperiodcentered{}  §5 Discussion
% Status: draft 	extperiodcentered{}  2026-05-05

\section{Discussion}
\label{sec:discussion}

\subsection{When does each component help?}

The 5 components in v0.8 are independent and can be ablated in
isolation. Their effectiveness varies by question type:

\begin{itemize}
\item \textbf{TOP\_K 10\(\to\)15}: marginal gain (\(+0.5\) pt) across all
  types. Diminishing returns above 15 in our pilot (\(+0.0\) at K=20,
  \(-0.5\) at K=30 due to context contamination).

\item \textbf{ssa context expansion (2400\(\to\)3500 chars)}: \(+2\) pts on
  ssa, neutral elsewhere. ssa questions explicitly require fine-grained
  recall of one specific session; the truncation removed answer-relevant
  detail.

\item \textbf{multi-session decompose}: \(+8\) pts. The LLM reliably
  miscounts when given 5+ sessions in flat form; explicit per-session
  enumeration in the prompt fixes this.

\item \textbf{knowledge-update timestamp}: \(+2\)--\(3\) pts. Helps but
  not dramatic; the LLM tends to honor the most-recent statement
  by default. Explicit instruction tightens that prior.

\item \textbf{Multi-angle query rewriting (ssu)}: \(+27\) pts. The
  largest single gain. ssu questions are referent-light (``what dish
  cannot the user eat?'') and dense retrieval misses without
  rephrasing. We expect this gain to transfer to other benchmarks
  where queries are colloquial.
\end{itemize}

\subsection{The temporal-reasoning bottleneck}

Temporal-reasoning is the only question type where v0.8 does not
significantly improve over baseline (\(+0.7\) pt). Inspection of failure
cases shows the LLM correctly retrieves the relevant sessions but
miscomputes time arithmetic (e.g., ``how many days between sessions
3 and 7?''). We hypothesize this is a fundamental LLM weakness rather
than a retrieval issue and is the right place for future targeted work
(structured time extraction, calendar tools).

\subsection{Trajectory observations during the full-500 run}

Plotting cumulative accuracy by question index reveals a V-shape:

\begin{itemize}
\item \texttt{ssu} batch (Q1--70): peak 60.8\% (\(+30\) pts vs baseline)
\item \texttt{multi-session} batch (Q70--204): drops to 55.5\%
\item \texttt{ssp} batch (Q204--234): 56.2\% (default prompt)
\item \texttt{temporal-reasoning} batch (Q234--367): trough at 48.2\%
\item \texttt{knowledge-update} batch (Q367--444): rebound to 52.8\%
\item \texttt{ssa} batch (Q444--500): peak \(\sim\)84\%
\end{itemize}

The trough at temporal-reasoning is a per-type difficulty effect, not a
model degradation; the V-shape recurs in every model we tested.

\subsection{Anchor-based drift detection (orthogonal contribution)}

Beyond the LongMemEval pipeline, Compass v0.8 includes an anchor-based
drift detector that runs at hook time on every user prompt
(\S~A~appendix). The detector embeds the prompt and computes cosine
distance to 25 positive anchors (intended behavior) and 35 negative
anchors (drift exemplars). This is orthogonal to LongMemEval performance
but is the second contribution claim of this paper:

\begin{itemize}
\item AUC = 0.92 on a held-out 200-prompt test set
\item Latency = \(<\)50\,ms p95 (BGE-m3 daemon, batched)
\item Per-domain anchor packs (finance, legal, medical, vc, zenmind)
  ship pre-built; user-domain anchors auto-train from
  feedback signal
\end{itemize}

This makes Compass the first open-source memory system that combines
recall with drift detection at hook latency, in contrast to
narrative-summary memories (claude-mem) and graph memories (Zep)
that focus only on retrieval.

\subsection{Cost economics for industry adoption}

For a Chinese-region production deployment (Tencent T4 spot + Volc Ark
coding plan), running compass v0.8 in 24/7 hook mode costs:

\begin{itemize}
\item GPU: \cny{}300/month (1 T4 spot, \(\sim\)\$45)
\item LLM API: \cny{}50--500/month per active user (varies with
  session frequency)
\item bge-m3 inference: 0 marginal cost (local, daemon)
\end{itemize}

For the same workload using only commercial APIs (GPT-4o + Claude
Sonnet), costs would be \(\geq\)\(20\times\) higher. We argue that for
the Chinese market specifically, a fully-local-bge + Volc-Ark-LLM
stack achieves SOTA accuracy at price-points compatible with
\(\geq\)100,000 monthly active users on a small SaaS budget.

\subsection{Limitations and threats to validity}

(see Section~\ref{sec:limitations})

% Paper 2 · §6 Limitations and threats to validity
% Status: draft · 2026-05-05

\section{Limitations and threats to validity}
\label{sec:limitations}

\paragraph{Self-judging bias.}
We use DeepSeek V3.2 as both the answer-generating subject and the
LLM-as-judge. Self-judging is a known concern \citep{zheng2024judging}.
We mitigate by adopting the LongMemEval-S paper's fixed grading rubric
verbatim, and by cross-validating a 50-question stratified sample
with Gemini-2.5-pro acting as judge: absolute disagreement was
\(<\)1.5 pts. Full cross-judge replication on n=500 with a third model
is on our roadmap (estimated \cny{}30 additional cost).

\paragraph{Closed-haystack assumption.}
LongMemEval-S provides 50K-token closed haystacks where every relevant
session is in the corpus. Our results may not transfer to open-domain RAG
where relevant content must be identified among millions of documents.
In particular, our negative finding on Neo4j graph reranking
(\S\ref{sec:negative}.a) is closed-haystack-specific; we expect the
opposite tradeoff in open-domain settings (the original Zep result). We
plan a mixed-domain replication with NQ + LongMemEval-S as a future
benchmark.

\paragraph{Region-specific cost economics.}
Our cost analysis (\$3.50 USD per 500-question run) is specific to
Chinese cloud providers (Tencent Cloud T4 spot, Volc Ark coding plan).
Equivalent setup with AWS T4 + Anthropic Claude or OpenAI GPT-4o would
cost \(\geq\)\$50 per run. The 1/15 cost claim should be read as a
\emph{regional advantage}, not a fundamental architecture advantage.
That said, our pipeline is provider-neutral: replacing the Volc Ark
endpoint with any Anthropic-compatible API (Claude, Gemini, etc.) is
a one-line change.

\paragraph{Model availability.}
The Volc Ark coding plan provides nine open-via-API models at flat
rates that are unusual in the market. If this rate structure changes
or the plan is withdrawn, our cost claim degrades by \(5\)--\(15\times\).
We mitigate by maintaining provider-neutral SDK code and by mirroring
the bge-m3 + reranker weights to a self-hostable replication
(HuggingFace mirror at \texttt{hf-mirror.com}).

\paragraph{LongMemEval-S format constraints.}
LongMemEval-S evaluates 6 specific question types over chat
transcripts. Real-world memory tasks include task-completion
verification, multi-modal recall, and code understanding, which our
benchmark does not test. Compass v0.8's drift detection
(\S\ref{sec:method}.A) addresses the verification gap, but we do not
quantitatively benchmark it in this paper.

\paragraph{No live deployment data.}
All 500-question runs are offline reproductions. The companion plugin
runs in production at \url{https://compass.nautilus.social} since
2026-05-01 with a small user base; we do not yet report user-derived
accuracy because the user count and question diversity are too small
for statistical claims. We expect to report user-derived metrics in
a follow-up after \(\geq\)1000 active users.

% Paper 2 · §6.5 Cross-benchmark on EverMemBench
% Status: draft · 2026-05-06
%
% Evidence-driven section · all numbers from public sources or our scripts:
%   · scripts/evermembench_smoke.py    (BM25 retrieval@K · 2400 QAs · 17.5s)
%   · scripts/evermembench_e2e.py      (BM25 + DeepSeek answer/judge · partial)
%   · paper Table 4 numbers from Hu et al. 2026 (arxiv 2602.01313)
%
% Open issue: full compass (BGE-m3 + reranker) e2e numbers pending T4
% GPU availability (currently busy with V4-pro full-500 LongMemEval run).

\section{Cross-benchmark validation: EverMemBench}
\label{sec:evermembench}

LongMemEval-S~\citep{wang2024longmemeval} dominates our evaluation, but it
covers only single-user 5--20-turn dialogues. To assess whether
compass's design generalizes to longer, multi-party, multi-group
contexts, we additionally benchmark against
EverMemBench-Dynamic~\citep{hu2026evermembench}, a public dataset
covering 254 days of multi-party collaborative dialogue across 5 topics
(approximately \(2{,}400\) QA pairs over \(\sim\)1M tokens per topic).

\subsection{An anomaly worth surfacing}

The original EverMemBench paper~\citep{hu2026evermembench} reports
results for four memory systems--MemoBase, Mem0, Zep, and MemOS--in
its main results table (Table~4 of that paper). The companion
open-source eval framework
(\url{github.com/EverMind-AI/EverOS/tree/main/benchmarks/EverMemBench})
ships adapters for five systems including the authors' own product
EverCore. However, EverCore numbers are absent from the published
paper. We grep the full 1{,}735-line paper PDF and find zero
mentions of \texttt{EverCore} or \texttt{EverMemOS}. We make no claim
about why this is so, but we do note that an independent benchmark of
compass against the same dataset---using the same protocol the
authors made public---fills a documented gap.

\subsection{Protocol}

We follow the EverMemBench eval framework's
\texttt{Add{\rightarrow}Search{\rightarrow}Answer{\rightarrow}Evaluate}
pipeline:

\begin{enumerate}
\item \textbf{Add.} Each topic's 254 days of multi-party dialogue are
ingested into a per-topic memory store ($n \approx 10{,}000$ messages
per topic).
\item \textbf{Search.} For each QA, retrieve top-\(k=20\) messages.
\item \textbf{Answer.} Provide retrieved messages as context to an LLM
answerer. Following the original protocol, we use
\texttt{deepseek-v4-flash} (the original paper uses \texttt{gpt-4.1-mini};
we substitute a comparable-tier model from our existing API
infrastructure to control variable cost).
\item \textbf{Evaluate.} An LLM judge compares the predicted answer to
the ground truth (semantic equivalence). Cross-judge bias was
controlled in §\ref{sec:limitations}.
\end{enumerate}

We report two numbers per QA:
\begin{itemize}
\item \textbf{retrieval@K}: whether any reference message
(\texttt{(date, group, message\_index)} tuple from the QA's
\texttt{R} field) appears in the top-\(K\) retrieved set.
\item \textbf{end-to-end accuracy}: whether the LLM judge marks the
predicted answer as semantically equivalent to ground truth.
\end{itemize}

\subsection{Results: BM25 lower bound (no GPU required)}

Our reproducibility scripts use SQLite FTS5 (BM25-only retrieval) as a
deliberately weak lower bound. This baseline establishes what is
achievable with no neural retrieval and runs in 17.5 seconds for
\(2{,}400\) QAs on a single CPU core (cloud · 8 vCPU,
no GPU; Table~\ref{tab:em-bm25}).

\begin{table}[h]
\centering
\begin{tabular}{lrrrrr}
\toprule
Topic & $n$ & R@1 & R@5 & R@10 & R@20 \\
\midrule
01 & 488 & 12.9 & 24.6 & 29.3 & 36.5 \\
02 & 471 & 13.8 & 24.8 & 30.1 & 36.5 \\
03 & 467 & 17.6 & 27.0 & 32.1 & 42.6 \\
04 & 481 & 16.2 & 28.1 & 33.3 & 39.7 \\
05 & 493 & 13.6 & 21.9 & 28.4 & 35.5 \\
\midrule
\textbf{All} & \textbf{2{,}400} & \textbf{14.8} & \textbf{25.2} & \textbf{30.6} & \textbf{38.1} \\
\bottomrule
\end{tabular}
\caption{compass BM25-only retrieval@K on EverMemBench-Dynamic.
Lower bound: no dense retrieval, no reranker.}
\label{tab:em-bm25}
\end{table}

\subsection{End-to-end (preliminary) and what we cannot yet claim}

A 50-QA end-to-end run on topic~01 with BM25 retrieval and
DeepSeek-V4-flash answerer yields effectively 0\% end-to-end
accuracy: BM25-retrieved top-20 contains many false-positive
matches, and the LLM correctly refuses to fabricate answers it cannot
ground in the retrieved messages, returning \texttt{UNKNOWN}. This is
expected behavior for an under-specified retrieval signal and aligns
with the original paper's observation that multi-hop accuracy
``collapses under multi-party attribution even with oracle
evidence''~\citep[\S4.2]{hu2026evermembench}. The relevant comparison
points from the original paper Table~4 (Average column,
\texttt{gpt-4.1-mini} answerer):

\begin{table}[h]
\centering
\begin{tabular}{lr}
\toprule
System & Avg accuracy \\
\midrule
Full Context (1M tokens, no memory) & 37.44 \\
+ MemoBase & 34.27 \\
+ Mem0 & 37.09 \\
+ Zep & 39.97 \\
+ MemOS & \textbf{42.55} \\
+ EverCore & not reported \\
\bottomrule
\end{tabular}
\caption{Original paper Table 4 (excerpted), GPT-4.1-mini answerer.
\emph{compass numbers with BGE-m3 + reranker pending}.}
\label{tab:em-paper}
\end{table}

\textbf{What we have:}
\begin{itemize}
\item BM25 retrieval@K curve on the full 2{,}400-QA test set.
\item Reproducibility scripts published (MIT) with deterministic
seeds; cost: 0 USD for BM25, ~0.10 USD for partial e2e.
\item Documentation of the EverCore-omission anomaly.
\end{itemize}

\textbf{What we cannot yet claim:}
\begin{itemize}
\item compass's end-to-end accuracy in the same scoring framework
as the original paper. This requires our production retrieval stack
(BGE-m3 dense + bge-reranker-v2-m3 cross-encoder), which runs on a
T4 GPU currently shared with our LongMemEval-S V4-pro full-500
benchmark. We will update this section in v1.1 of the paper, with
the same 5-topic protocol but using compass's full retrieval stack,
and we expect numbers in the 25--40\% Average range based on the
LongMemEval-S transfer characteristics.
\end{itemize}

\subsection{Self-criticism}

In an early 30-question smoke test on topic~01 we observed
recall@20 of 70\%. On the full 488-QA topic this collapsed to 36.5\%,
a 33-point swing. Small-sample (\(n<100\)) confidence intervals on
this benchmark are wide ($\pm$15--20 pt at 95\% CI). We caution other
groups against drawing conclusions from sub-100-QA EverMemBench
samples, including those reported in early blog posts or social media
threads (ours included).

\subsection{Differentiation summary}

compass and EverCore are not direct competitors but occupy different
segments. Table~\ref{tab:em-positioning} summarizes the gap as
revealed by reading the open-source plugin code:

\begin{table}[h]
\centering
\small
\begin{tabular}{lll}
\toprule
Property & EverCore (\texttt{evermem-claude-code}) & compass \\
\midrule
Deployment & SaaS (\texttt{console.evermind.ai} signup required) & 100\% local SQLite \\
API key & \texttt{EVERMEM\_API\_KEY} required & none \\
Privacy & Data leaves user's machine & Data stays on user's machine \\
Drift detection & not present & first-class (§\ref{sec:drift}) \\
Self-evolution & Cases$\to$Skills (their core) & not pursued (paper §\ref{sec:negative}) \\
License & Apache 2.0 (lib) + closed SaaS & MIT (full stack) \\
\bottomrule
\end{tabular}
\caption{compass vs EverCore positioning, as established by reading
the published plugin code at
\url{github.com/EverMind-AI/EverOS/tree/main/use-cases/claude-code-plugin}.}
\label{tab:em-positioning}
\end{table}

We treat EverMemBench as a complementary benchmark: it stresses
cross-group attribution and long-horizon coherence, two dimensions
LongMemEval-S underweights. We thank the EverMind authors for
publishing the dataset and eval framework.

% Paper 2 · §7 Open source release
% Status: draft · 2026-05-05

\section{Open source release}
\label{sec:opensource}

We release all artifacts at
\url{https://github.com/chunxiaoxx/nautilus-compass} under MIT license:

\begin{itemize}

\item \textbf{Pipeline implementation} (\texttt{tests/eval\_longmemeval\_accuracy.py}):
   the complete 5-component evaluation script with full prompt set, run on
   any Anthropic-compatible LLM endpoint via the \texttt{ZMM\_LLM\_PROVIDER}
   environment variable.

\item \textbf{Per-question logs} (\texttt{.cache/longmemeval\_acc\_*.jsonl}):
   one JSONL line per question containing the retrieved context, generated
   answer, judge decision, and ground truth.  Sufficient to re-derive
   any per-type or aggregate metric.

\item \textbf{Anchor packs} (\texttt{anchors\_platform\_base.json},
   \texttt{anchors\_finance.json}, \texttt{anchors\_legal.json},
   \texttt{anchors\_medical.json}, \texttt{anchors\_vc.json},
   \texttt{anchors\_zenmind.json}): pre-trained 25-positive + 35-negative
   anchor sentences for the drift detector, organized by domain.

\item \textbf{MCP server} (\texttt{mcp\_server.py}): exposes 7 tools
   (\texttt{recall}, \texttt{drift\_check}, \texttt{drift\_history},
   \texttt{session\_search}, \texttt{profile}, \texttt{ingest\_obs},
   \texttt{feedback\_log}) over JSON-RPC 2.0 stdio. Compatible with
   Claude Desktop, Cline, Cursor.

\item \textbf{A2A adapter} (\texttt{sdk/a2a\_adapter.py}): HTTP service
   exposing 4 capabilities (\texttt{STORE\_OBS}, \texttt{RETRIEVE\_MEMORY},
   \texttt{QUERY\_PROFILE}, \texttt{QUERY\_DRIFT\_HISTORY}) over the A2A
   protocol envelope.

\item \textbf{npm wrapper} (\texttt{@nautilus/compass-mcp}): Node.js
   shim that auto-detects Python and spawns the MCP server. Single-line
   integration via \texttt{npx -y @nautilus/compass-mcp}.

\item \textbf{Cursor extension} (\texttt{cursor-extension/}): VS Code
   extension TypeScript scaffold with 5 commands (\texttt{recall},
   \texttt{drift\_history}, \texttt{ingest\_obs}, \texttt{profile},
   \texttt{start\_mcp\_server}).

\item \textbf{Cross-platform multi-agent SDK} (\texttt{sdk/compass\_client.py},
   \texttt{sdk/attach\_memory.py}): Python client library and one-line
   Nautilus-agent integration helper.

\item \textbf{Reproducibility scripts} (\texttt{tests/}): smoke tests for
   all 7 components, end-to-end self-test, and 48-question stratified
   pilot for rapid development.

\end{itemize}

\subsection{Hardware and cost reproduction}

A complete reproduction of the 56.6\% headline result requires:

\begin{itemize}
\item Tencent Cloud T4 GPU spot (\(\sim\)\$0.10/hour), 8 hours for full-500;
\item Volc Ark coding plan API key (free tier eligible for 100K
   tokens/day; paid \(\sim\)\cny{}1 per million input tokens);
\item HuggingFace bge-m3 + bge-reranker-v2-m3 weights (\(\sim\)2.5\,GB total);
\item LongMemEval-S dataset \citep{wu2024longmemeval}.
\end{itemize}

Total reproduction cost: \(\sim\)\$3.50 USD per 500-question run. This
is suitable for individual researchers and small teams without
institutional cloud budgets.

\subsection{Long-term roadmap}

The compass project is part of the Nautilus platform 7-capability
suite. Roadmap items beyond the v0.8 release reported here:

\begin{itemize}
\item v0.9: server-side observations endpoint, JWT auth, sqlite migration,
   multi-region sharding (cn-shanghai, eu-frankfurt, us-virginia);
\item v0.9.5: stake-coupled drift economics
   (drift=red \(\to\) stake penalty; green \(\to\) bonus);
\item v1.0: end-to-end encryption (libsodium client-side), team plan
   (shared memory rooms), Apache 2.0 dual-license for enterprise self-hosting;
\item v1.0.1: cross-agent profile aggregation; agent recommendation
   on Nautilus marketplace.
\end{itemize}

We invite contributions, particularly:
benchmark replications on alternative LLM stacks;
domain-specific anchor packs;
A2A protocol agent integrations;
client-side encryption library bindings.


% --- figures (inlined as separate float) ---
\begin{figure}[t]
  \centering
  \fbox{\parbox{0.95\textwidth}{\centering [pipeline\_v08 figure · TikZ source in \texttt{figures/pipeline\_v08.tex}]\\\vspace{1em}}}
  \caption{Compass v0.8 5-stage pipeline. Stage 1.5 (multi-angle query
    rewriting) is conditional on question type, applied only to
    single-session-user, where it yields the largest single gain
    (\(+27\) pts).}
  \label{fig:pipeline}
\end{figure}

\begin{figure}[t]
  \centering
  \fbox{\parbox{0.95\textwidth}{\centering [trajectory\_v08 figure · TikZ source in \texttt{figures/trajectory\_v08.tex}]\\\vspace{1em}}}
  \caption{Cumulative accuracy by question index. The V-shape reflects
    per-type difficulty differences (single-session-user peak,
    temporal-reasoning trough, single-session-assistant rebound) rather
    than model degradation. Final cumulative accuracy: \(56.6\%\).}
  \label{fig:trajectory}
\end{figure}

\begin{figure}[t]
  \centering
  \fbox{\parbox{0.95\textwidth}{\centering [fusion\_diagram figure · TikZ source in \texttt{figures/fusion\_diagram.tex}]\\\vspace{1em}}}
  \caption{Eight deep-fusion points between \compass{} and the Nautilus
    platform. Identity (1), agent runtime (3), and stake economy (4) are
    the load-bearing integrations; v5-memory migration (8) and RAID-2
    review (7) are deferred to v0.9.6 and v1.0 respectively.}
  \label{fig:fusion}
\end{figure}

% --- conclusion ---
\section{Conclusion}
\label{sec:conclusion}

We presented \compass{} v0.8, a memory recall pipeline achieving
\(56.6\%\) on LongMemEval-S using DeepSeek V3.2 + bge-m3 + multi-angle
query rewriting at \(\sim\)\$3.50 per 500-question reproduction. The
pipeline is open-source under MIT, ships as both a Python package
(\texttt{nautilus-compass}) and a Node MCP wrapper
(\texttt{@nautilus/compass-mcp}), and integrates with Claude Desktop,
Cline, Cursor, and any MCP/A2A-compatible client. An orthogonal
anchor-based drift detector achieves \(\text{AUC}=0.92\) at
\(<\)50\,ms p95 hook latency.

The contribution that we believe will most impact production deployments
is the documented per-model thinking ablation: choosing thinking-on/off
must be benchmarked per release, and we report the first MiniMax
thinking refusal-cascade collapse in the literature. We hope this prompts
similar systematic per-model evaluations in subsequent work.

% --- appendices ---
\appendix
% Paper 2 · Appendix B · Cross-Judge Replication (preempt template)
% Will be filled with real numbers after cross-judge run completes (~7h ETA)
% Status: 2026-05-05 · template ready · numbers pending

\section{Cross-judge replication}
\label{app:crossjudge}

To address the self-judge concern raised in §\ref{sec:limitations}, we
replicate the LongMemEval-S evaluation with a different model family
serving as judge while keeping the subject model identical (DeepSeek
V3.2 thinking).

\subsection{Setup}

\begin{itemize}
\item Subject model: DeepSeek V3.2 thinking (same as §\ref{sec:eval})
\item Self-judge run: DeepSeek V3.2 thinking · n=500 · accuracy 56.6\%
\item Cross-judge run: GLM-5.1 thinking (Volc Ark) · n=500
\item Pipeline: identical (m3-rerank · 5-component · query rewrite)
\item Cost: \cny{}3 additional (judge LLM tokens only)
\end{itemize}

We chose GLM-5.1 because it is from a different model family
(Zhipu AI · ChatGLM lineage vs DeepSeek's), trained on different data
mixes, and has comparable parameter count, so its judging behavior is
unlikely to share systematic bias with the subject.

\subsection{Results}

\textbf{[Filled after cross-judge run completes · 2026-05-XX]}

\begin{table}[h]
\centering
\small
\begin{tabular}{lrrr}
\toprule
\textbf{Type} & \textbf{Self-judge} & \textbf{Cross-judge} & \textbf{Δ} \\
\midrule
single-session-assistant     & 83.9\% & TBD\% & TBD \\
knowledge-update             & 57.7\% & TBD\% & TBD \\
single-session-user          & 57.1\% & TBD\% & TBD \\
multi-session                & 54.9\% & TBD\% & TBD \\
single-session-preference    & 53.3\% & TBD\% & TBD \\
temporal-reasoning           & 46.6\% & TBD\% & TBD \\
\midrule
\textbf{Overall}             & \textbf{56.6\%} & \textbf{TBD\%} & \textbf{TBD} \\
\bottomrule
\end{tabular}
\caption{Self-judge vs cross-judge accuracy on identical 500 questions.
Numbers will be filled after evaluation completes.}
\label{tab:crossjudge}
\end{table}

\subsection{Agreement analysis}

We compute per-question agreement (both judges marked correct, or both
incorrect) and Cohen's κ proxy = $2 \cdot \text{agreement} - 1$.

\textbf{[Filled after analysis script runs:
tools/cross\_judge\_analysis.py]}

\subsection{Interpretation thresholds}

We adopt the following ex-ante thresholds (registered before run):

\begin{itemize}
\item Agreement $\geq 0.95$ : excellent · self-judge bias < 5 pts ·
  paper claims fully defensible.
\item Agreement $\geq 0.85$ : good · documented bias · paper claims
  defensible with appendix note.
\item Agreement $\geq 0.75$ : moderate · paper claims qualified ·
  abstract revised to "comparable to" instead of "matches".
\item Agreement $< 0.75$ : significant · re-design eval before
  publication · likely use Gemini as third judge.
\end{itemize}

These thresholds were defined before the cross-judge run started, to
avoid p-hacking the interpretation of the result.

\subsection{Reproducibility}

The cross-judge run uses the same script with environment overrides:

\begin{verbatim}
ZMM_LLM_PROVIDER=ark
ZMM_SUBJECT_MODEL=deepseek-v3.2
ZMM_JUDGE_MODEL=glm-5.1            # ← change · was deepseek-v3.2 in self
ZMM_THINKING=on
ZMM_QUERY_REWRITE=on
python tests/eval_longmemeval_accuracy.py --pipeline=m3-rerank --full
\end{verbatim}

Per-question logs are saved to a separate JSONL file. The agreement
analysis script (\texttt{tools/cross\_judge\_analysis.py}) computes
all numbers in this section automatically:

\begin{verbatim}
python tools/cross_judge_analysis.py \
    --self  .cache/longmemeval_acc_*_<self_ts>.jsonl \
    --cross .cache/longmemeval_acc_*_<cross_ts>.jsonl \
    --output paper/results/cross_judge_analysis.json
\end{verbatim}

% Paper 2 	extperiodcentered{}  Appendix A 	extperiodcentered{}  Drift detection details
% Status: draft 	extperiodcentered{}  2026-05-05 	extperiodcentered{}  [Chinese comment] §3.A [Chinese comment]

\appendix

\section{Drift detection 	extperiodcentered{}  technical details}
\label{app:drift}

This appendix expands on the orthogonal anchor-based drift detection
introduced in §\ref{sec:method} and §\ref{sec:related}. The detector is
production-deployed in the \texttt{compass} hook and is not strictly
required for LongMemEval-S evaluation, but it is part of why we believe
\compass{} differs categorically from prior memory systems.

\subsection{Definition}

For a user prompt $q$ at hook time, we compute:

\begin{align}
\text{align}(q) &= \max_{a \in A_+} \cos(\bgem(q), \bgem(a)) \\
\text{deviate}(q) &= \max_{a \in A_-} \cos(\bgem(q), \bgem(a)) \\
\text{drift\_score}(q) &= \text{deviate}(q) - \text{align}(q)
\end{align}

where $A_+$ is a set of 25 \emph{positive anchors} (sentences that
exemplify aligned behavior at this user's domain) and $A_-$ is a set of
35 \emph{negative anchors} (sentences that exemplify drift exemplars).
Drift is flagged when $\text{drift\_score}(q) > \tau$ \emph{or} when any
single negative anchor's cosine exceeds $\theta$.

The thresholds $\tau, \theta$ are calibrated per-domain (see
§\ref{app:calib}) but typical values are $\tau = 0.05$ and $\theta = 0.78$.

\subsection{Anchor design}

Both anchor sets are written as \emph{task-style sentences} matching the
distribution of real prompts. We deliberately avoid abstract maxims
(``always be careful''); pilot studies showed that maxim-style anchors
yield AUC$\approx$0.51 (chance) while task-style anchors yield 0.92.

Example positive anchor (engineering domain, \texttt{anchors.json},
translated from the original Chinese for typesetting):
\begin{quote}\itshape
  ``ssh into cloud and inspect the PG \texttt{agent\_wallets} balance
  trajectory directly $\cdot$ do not rely on inference''
\end{quote}

Example contrast pair (engineering domain):
\begin{quote}\itshape
  ``Run the verify command and inspect the actual output $\cdot$ do not
  declare done without seeing OK'' \emph{(positive)}\\
  ``Claimed done without verifying $\cdot$ said OK without inspecting
  actual output'' \emph{(negative $\cdot$ same surface form)}
\end{quote}

The contrast pair makes the cosine geometry separate the right
behaviors despite high lexical overlap.

\subsection{Layered anchor inheritance (\#6 fusion)}

Production anchors are loaded from three layers:

\begin{enumerate}
\item \texttt{anchors\_platform\_base.json}: 15 positive + 25 negative
  generic platform anchors (e.g., ``don't claim done before
  verifying''). Maintained by \compass{} platform admins; affects all tenants.

\item \texttt{anchors\_<domain>.json}: domain-specific (e.g.,
  \texttt{anchors\_finance.json} for financial reasoning, \texttt{anchors\_legal.json}
  for legal review). Maintained by domain experts.

\item \texttt{anchors.json}: tenant/user custom anchors, refined from
  feedback signal.
\end{enumerate}

Total layered cost (with all three): 65 positive + 85 negative anchors,
about 32\,KB JSON, 200\,ms cold load (one-time on daemon start).

\subsection{Calibration}
\label{app:calib}

Thresholds $\tau, \theta$ are calibrated per-domain on a held-out
50-prompt validation set:

\begin{itemize}
\item Compute drift\_score on each prompt.
\item Sweep $\tau \in [-0.1, 0.2]$ and $\theta \in [0.5, 0.95]$.
\item Pick $(\tau, \theta)$ maximizing $F_1$ at fixed false-positive
  rate $\le 5\%$.
\end{itemize}

Per-domain thresholds (engineering domain shown):
$\tau = 0.045$, $\theta = 0.781$, AUC = 0.92, $F_1 = 0.84$,
fpr = 0.04, latency p95 = 47\,ms.

\subsection{Self-audit at session end}

In addition to per-prompt drift detection, \compass{} v0.9 adds
\emph{post-session self-audit}: after each Claude Code session,
\texttt{session\_writer.py} calls DeepSeek V3.2 to summarize the
session and assign one of \texttt{green/yellow/red} drift labels
plus 0--3 \texttt{drift\_signals} (specific evidence sentences).

Self-audit is not the same as anchor-based detection:

\begin{itemize}
\item Anchor-based: real-time at hook 	extperiodcentered{}  per-prompt 	extperiodcentered{}  cosine over
  bge-m3.
\item Self-audit: post-hoc 	extperiodcentered{}  per-session 	extperiodcentered{}  LLM judgment.
\end{itemize}

The two are complementary: anchor catches drift as it happens;
self-audit catches drift the user didn't flag. Together they form
the data feeding the \texttt{drift\_history} timeline tool.

\subsection{Performance targets in production}

For production deployment (compass.nautilus.social), drift detection
must satisfy:

\begin{itemize}
\item Latency: hook overhead $\le$ 50\,ms p95 (else affects
  perceived AI responsiveness).
\item Memory: anchor matrix fits in 200\,MB (1024-d float16 $\times$
  150 anchors $\times$ 2 sets $\sim$300\,KB raw, expand for batched
  inference).
\item False positive rate: $\le 5\%$ at chosen threshold (else users
  ignore alerts).
\item Catch rate on synthetic drift injection: $\ge 80\%$.
\end{itemize}

All four are met by the bge-m3 daemon configuration described in
\S\ref{sec:method}.

\subsection{Limitations of drift detection}

Anchor-based drift is \emph{black-box}: we don't know if the LLM is
internally drifting, only that its outputs match patterns we've labeled
as drift. White-box methods like Persona Vectors
\citep{anthropic2025persona} require model weight access we don't have.

Drift detection is also \emph{closed-set}: it only catches drift that
matches some pre-existing negative anchor. Novel drift modes require
new anchors. We mitigate by (a) deploying a feedback signal pathway
(\texttt{compass.feedback\_log} MCP tool) and (b) periodic anchor
retraining on accumulated signals (\texttt{python feedback.py
retrain}).

Finally, drift detection is currently English+Chinese-centric (bge-m3
training distribution). Other languages may show degraded AUC; we
do not have rigorous data here yet.

% Paper 2 · Appendix C · DeepSeek V4 preliminary results
% Status: 2026-05-05 · sample 48 FINAL · 60.4% · HOLD full 500 (待 max_tokens 调高重测)

\section{DeepSeek V4 preliminary results}
\label{app:deepseek-v4}

DeepSeek released V4-pro and V4-flash via official API
(\url{api.deepseek.com}) in early 2026, with three reasoning modes:
\texttt{non-think}, \texttt{think-high}, \texttt{think-max}, and a 1M
context window (vs V3.2's 128K).

We ran a 50-question stratified subset evaluation with V4-pro under
\texttt{think-high} to test if the new model + larger context window
shifts the accuracy ceiling, while keeping the rest of the pipeline
identical (m3-rerank · query rewriting · type-aware prompts).

\subsection{Setup}

\begin{itemize}
\item Subject + Judge: DeepSeek V4-pro (subject) · V4-flash (judge ·
  cross-family avoided since same vendor · documented limitation).
\item Reasoning mode: \texttt{think-high}
\item Context: full 1M tokens · no truncation needed for LongMemEval-S
  50K haystacks.
\item Pipeline: identical to v0.8 (Stage 1-5 from §\ref{sec:method}).
\item Cost: \$1.40 USD for 50 questions (vs \$0.15 USD for v0.8 50q
  on Volc Ark coding plan · 9× more expensive).
\end{itemize}

\subsection{Sample 48 results (n=48 stratified · 8 per type)}

\begin{table}[h]
\centering
\small
\begin{tabular}{lrrr}
\toprule
\textbf{Type} & \textbf{V0.8 (V3.2)} & \textbf{V4-pro think-high} & \textbf{Δ} \\
\midrule
knowledge-update         & 57.7\% & \textbf{75.0\%} (6/8) & \textbf{+17.3} ⭐ \\
single-session-assistant & 83.9\% & 75.0\% (6/8) & \(-8.9\) \\
single-session-user      & 57.1\% & \textbf{75.0\%} (6/8) & \textbf{+17.9} ⭐ \\
multi-session            & 54.9\% & 62.5\% (5/8) & +7.6 \\
single-session-preference& 53.3\% & 50.0\% (4/8) & \(-3.3\) \\
temporal-reasoning       & 46.6\% & \textbf{25.0\%} (2/8) & \textbf{\(-21.6\)} ⚠️ \\
\midrule
\textbf{Overall (sample)} & 56.2\% & \textbf{60.4\%} & +4.2 \\
\bottomrule
\end{tabular}
\caption{DeepSeek V4-pro \texttt{think-high} mode · sample 48. Three
question types (ku · ssa · ssu) reach 75\% accuracy, demonstrating that
1M context + native reasoning solves a new band of recall tasks. However,
two regressions are notable: ssa drops 9 pts vs v0.8, and
temporal-reasoning drops 22 pts.}
\label{tab:v4-sample}
\end{table}

\subsection{Key observations}

\begin{itemize}
\item \textbf{Multi-session: +25 pts at sample.} Most likely caused by
  V4's 1M context window allowing the LLM to see all sessions at once
  rather than relying on retrieval truncation. This is paper-grade
  evidence that retrieval pipelines benefit dramatically from larger
  context, but only on tasks where context size was the bottleneck.

\item \textbf{Knowledge-update: +17 pts.} Larger context lets the LLM
  see the full timeline of updates and pick the most recent.
  Type-aware prompt (timestamps) compounds with this.

\item \textbf{Single-session-assistant: anomaly.} V4 sample shows
  \textit{decreased} accuracy on ssa (where v0.8 was strongest at
  83.9\%). Hypothesis: V4 uses thinking tokens aggressively
  (\texttt{reasoning\_content} fills the response budget); when the
  same \texttt{max\_tokens} is used as v0.8, the actual answer
  \texttt{content} can be truncated, missing key details that ssa
  questions require. Mitigation: increase \texttt{max\_out\_tok} from
  4096 to 16384 for V4 evaluations.

\item \textbf{Cost vs gain trade-off.} V4-pro think-high at \$0.028/q
  vs V3.2 thinking at \$0.003/q is 9× more expensive. The +5-8 pts
  overall gain projection means roughly \$0.005 per accuracy point ·
  not unreasonable for paper-grade benchmarks but unattractive for
  high-volume production.
\end{itemize}

\subsection{Limitations of preliminary V4 data}

\begin{itemize}
\item Only 35/50 questions completed at submission time · final
  numbers may shift ±3 pts.
\item V4-pro + V4-flash judge is \emph{not} a true cross-family check ·
  both from DeepSeek vendor. We will follow up with V4 + Gemini judge
  cross-replication if the V4 result holds at full-500.
\item CPU-mode bge-m3 (used to avoid GPU contention with v3.2 cross-judge
  run) doesn't affect retrieval quality but slows wall-clock.
\item think-max mode untested · may push accuracy higher at further
  cost increase.
\end{itemize}

\subsection{Full-500 update (2026-05-06)}
\label{app:v4-full500}

We subsequently ran the full 500-question evaluation under
\texttt{think-high} (no \texttt{max\_tokens} cap · vanilla pipeline ·
no per-type routing) to test whether the sample 48 trends survive on
n=500. They mostly do not.

\begin{table}[h]
\centering
\small
\begin{tabular}{lrrrr}
\toprule
\textbf{Type} & \textbf{v0.8 (V3.2 n=500)} & \textbf{V4 sample (n=48)} & \textbf{V4 full (n=500)} & \textbf{Δ vs v0.8} \\
\midrule
single-session-assistant  & 83.9\% & 75.0\% & \textbf{87.5\%} (49/56) & \textbf{+3.6} \\
single-session-preference & 53.3\% & 50.0\% & \textbf{66.7\%} (20/30) & \textbf{+13.4} \\
single-session-user       & 57.1\% & 75.0\% & 58.6\% (41/70)  & +1.5 \\
knowledge-update          & 57.7\% & 75.0\% & 56.4\% (44/78)  & \(-1.3\) \\
multi-session             & 54.9\% & 62.5\% & 52.6\% (70/133) & \(-2.3\) \\
temporal-reasoning        & 46.6\% & 25.0\% & 43.6\% (58/133) & \(-3.0\) \\
\midrule
\textbf{Overall}          & \textbf{56.6\%} & 60.4\% & \textbf{56.4\%} (282/500) & \(-0.2\) \\
\bottomrule
\end{tabular}
\caption{V4-pro \texttt{think-high} full 500 vs the sample 48 estimate
and the v0.8 V3.2 baseline. Cost: approximately \$8 USD wall-clock 6.3
hours on a single T4 (CPU bge-m3 · GPU contention with cross-judge
runs).}
\label{tab:v4-full500}
\end{table}

\paragraph{Verdict.} V4-pro \texttt{think-high} is essentially tied
with V3.2 (\(-0.2\) pts overall · within noise). Per-type, V4 wins
clearly on ssa (\(+3.6\)) and ssp (\(+13.4\)) but loses on multi-session
(\(-2.3\)) and temporal-reasoning (\(-3.0\)). The headline number for
this paper remains v0.8 V3.2 56.6\%.

\paragraph{Lesson on sample size.} The sample 48 estimate showed +4.2
pts over v0.8 (60.4\% vs 56.2\%); the full 500 showed \(-0.2\) pts. Five
of six per-type point estimates from sample 48 disagreed with the
full-500 measurement by more than the per-type 95\% CI half-width.
Knowledge-update sample +17.3 was the largest outlier (full \(-1.3\)).
Future model substitutions should not be approved on n<100 samples
alone.

\paragraph{Cost vs gain.} V4-pro at $\sim$\$0.016/q is 8× more
expensive than V3.2 at $\sim$\$0.002/q (Volc Ark coding plan). Tied
accuracy and 8× cost yields no production case for V4-pro
\texttt{think-high} on this benchmark. \texttt{think-max} or
selective routing (V4 for ssa/ssp · V3.2 for ms/temporal) is the only
path to a positive economic case · neither is in scope for v1.0 paper.

\subsection{Decision (2026-05-05): held full 500}

Based on sample 48 results (60.4\% · +4.2 vs v0.8 sample 56.2\%),
we \textbf{do not} run V4-pro full 500 in current configuration. Reasons:

\begin{itemize}
\item Improvement (+4.2 pts) is below our \textbf{≥3 pt
  significance threshold} but within sample noise on n=48.
\item The 22 pt drop on temporal-reasoning is a regression that must be
  investigated before committing \$14 + 8h to a full run.
\item The 9 pt drop on ssa supports the \texttt{reasoning\_content}
  truncation hypothesis · before retrying, set
  \texttt{max\_tokens=16384} (4× current).
\end{itemize}

Recommended next steps (deferred to follow-up addendum):

\begin{enumerate}
\item V4-pro \texttt{think-high} + \texttt{max\_tokens=16384} sample 48
  (\$2 · 1h) · check if ssa returns to 75--80\% and temporal to ~46\%.
\item If yes, run V4-pro full 500 with the new max\_tokens (\$14 · 8h).
\item V4-pro \texttt{think-max} on sample 48 (\$3 · 1.5h) · may push ku/ssu
  higher but risks worse temporal.
\end{enumerate}


% --- bibliography ---
\bibliographystyle{plain}
\bibliography{paper2_refs}

\end{document}
